\begingroup
\raggedbottom
\setlength{\emergencystretch}{2em}
\section{概形的极限}\label{src:B.3}
在模理论中, 我们经常需要处理非 Noether 概形, 原因很简单: 模函子和叠定义在所有概形构成的范畴上. 若改为在局部 Noether 概形的范畴中工作, 则有一个局限: 此范畴对纤维积不封闭; 而域或 $\mathbb Z$ 上有限型概形的范畴不包含概形的局部环或其完备化. 无论如何, 使用本节介绍的极限方法, 通常很容易将概形及其态射的性质归结为 Noether 情形.

\subsection{极限的存在性与 Noether 逼近}\label{src:B.3.1-subsection}
过渡映射仿射的概形逆系统存在极限.
\begin{proposition}[极限的存在性]\label{src:B.3.1}
若 $(S_\lambda,f_{\lambda\mu})_{\lambda\in\Lambda}$ 为过渡映射仿射的概形逆系统, 则在概形范畴中存在极限 $S=\varprojlim_{\lambda\in\Lambda}S_\lambda$, 且每个态射 $f_\lambda:S\to S_\lambda$ 均仿射.
\end{proposition}
\begin{proof}
若每个 $S_\lambda=\Spec A_\lambda$ 都仿射, 则取 $S=\Spec(\varinjlim_\lambda A_\lambda)$. 一般情形下, 选取元素 $0\in\Lambda$, 并令
\[
S=\Spec_{S_0}\left(\varinjlim_{\lambda\geq0}f_{\lambda0,*}\mathcal O_{S_\lambda}\right).
\]
细节见 \cite[IV.$8.2$]{EGA} 及 \cite[Tag 01YX]{SP}.
\end{proof}
每个仿射概形 $\Spec A$ 都是 $\mathbb Z$ 上有限型仿射概形 $\Spec A_\lambda$ 的极限. 这由环 $A$ 是其有限生成 $\mathbb Z$-子代数之并这一事实推出. 更一般地, 有:
\begin{proposition}[相对 Noether 逼近]\label{src:B.3.2}
\index{relative Noetherian approximation@相对 Noether 逼近\quad relative Noetherian approximation|hyperpage}
设 $X\to S$ 为概形态射, $X$ 拟紧且拟分离, $S$ 拟分离. 则 $X=\varprojlim_{\lambda\in\Lambda}X_\lambda$ 是一个逆系统 $(X_\lambda,f_{\lambda\mu})$ 的极限, 其中 $X_\lambda$ 是 $S$ 上的有限呈示概形, 过渡映射为 $S$ 上的仿射态射.
\end{proposition}
\begin{proof}
见 \cite[Tag 09MV]{SP}. 当 $S=\Spec\mathbb Z$ 时, 此结论通常称为\newterm{绝对 Noether 逼近}{Absolute Noetherian Approximation}, 最早证明于 \cite[定理 C.$9$]{TT90}.
\end{proof}
\begin{proposition}[极限下态射的下降]\label{src:B.3.3}
设 $S=\varprojlim_{\lambda\in\Lambda}S_\lambda$ 为拟紧且拟分离概形逆系统的极限, 其过渡映射仿射.
\begin{enumerate}[(1)]
\item 对于概形之间的有限呈示态射 $X\to S$, 存在指标 $0\in\Lambda$ 和概形之间的有限呈示态射 $X_0\to S_0$, 使得 $X\simeq X_0\times_{S_0}S$. 此外, 若对 $\lambda>0$ 定义 $X_\lambda:=X_0\times_{S_0}S_\lambda$, 则 $X=\varprojlim_{\lambda\geq0}X_\lambda$ 为逆系统 $(X_\lambda,f_{\lambda\mu})$ 的极限, 其中对 $\lambda\geq\mu$, (仿射) 过渡映射 $f_{\lambda\mu}:X_\lambda\to X_\mu$ 为 $S_\lambda\to S_\mu$ 的基变换.
\item 对某个指标 $0\in\Lambda$, 设 $X_0$ 和 $Y_0$ 为 $S_0$ 上的有限呈示概形. 对 $\lambda>0$, 令 $X_\lambda=X_0\times_{S_0}S_\lambda$ 和 $Y_\lambda=Y_0\times_{S_0}S_\lambda$, 并设 $X=\varprojlim_\lambda X_\lambda$ 和 $Y=\varprojlim_\lambda Y_\lambda$ 为其极限 (命题 \bookref{B.3.1}). 则自然映射
\[
\varinjlim_{\lambda\geq0}\operatorname{Mor}_{S_\lambda}(X_\lambda,Y_\lambda)
\longrightarrow\operatorname{Mor}_S(X,Y)
\]
为双射.
\end{enumerate}
\end{proposition}
换言之, $S$ 上有限呈示概形的范畴是 $S_\lambda$ 上有限呈示概形范畴的余极限.
\begin{proof}
见 \cite[IV.$8.8$]{EGA} 及 \cite[Tag 01ZM]{SP}.
\end{proof}
\newpage
拟凝聚层也可以在极限下下降.
\begin{proposition}[极限下拟凝聚层的下降]\label{src:B.3.4}
设 $(S_\lambda,f_{\lambda\mu})$ 为拟紧且拟分离概形的逆系统, 过渡映射仿射, 极限为 $S=\varprojlim_{\lambda\in\Lambda}S_\lambda$. 记投影映射为 $f_\lambda:S\to S_\lambda$.
\begin{enumerate}[(1)]
\item 若 $F$ 为有限呈示拟凝聚 $\mathcal O_S$-模 (相应地, 向量丛, 线丛), 则存在指标 $\lambda\in\Lambda$ 和有限呈示 $\mathcal O_{S_\lambda}$-模 $F_\lambda$ (相应地, 向量丛, 线丛), 使得 $F\simeq f_\lambda^*F_\lambda$.
\item 对指标 $0\in\Lambda$, 设 $F_0$ 和 $G_0$ 为有限呈示 $\mathcal O_{S_0}$-模, $F$ 和 $G$ 为沿 $f_0$ 到 $S$ 的拉回, $F_\lambda$ 和 $G_\lambda$ 为沿 $f_{0\lambda}$ 到 $S_\lambda$ 的拉回. 自然映射
\[
\varinjlim_{\lambda\geq0}\operatorname{Hom}_{\mathcal O_{S_\lambda}}(F_\lambda,G_\lambda)
\longrightarrow\operatorname{Hom}_{\mathcal O_S}(F,G)
\]
为双射.
\item 对指标 $0\in\Lambda$, 设 $f_0:X_0\to Y_0$ 为 $S_0$ 上概形之间的有限呈示态射, $F_0$ 为 $X_0$ 上的有限呈示拟凝聚层. 若 $F_0$ 沿 $X_0\times_{S_0}S\to X_0$ 的拉回在 $Y_0\times_{S_0}S$ 上平坦, 则对 $\lambda\gg0$, $F_0$ 沿 $X_0\times_{S_0}S_\lambda\to X_0$ 的拉回在 $Y_0\times_{S_0}S_\lambda$ 上平坦.
\end{enumerate}
\end{proposition}
换言之, $S$ 上有限呈示模的范畴是 $S_\lambda$ 上有限呈示模范畴的余极限. 注意, 在 (2) 中取 $F_0=\mathcal O_{S_0}$, 可得 $\Gamma(S,F)=\varinjlim_{\lambda\geq0}\Gamma(S_\lambda,F_\lambda)$.
\begin{proof}
见 \cite[IV.$8.5.2$]{EGA} 及 \cite[Tags 01ZR, 0B8W 及 05LY]{SP}.
\end{proof}

\subsection{极限下性质的下降}\label{src:B.3.2-subsection}
\begin{proposition}[极限下概形性质的下降]\label{src:B.3.5}
设 $S=\varprojlim_\lambda S_\lambda$ 为拟紧且拟分离概形逆系统的极限, 其过渡映射仿射. 若 $S$ 仿射 (相应地, 拟仿射, 分离), 则对 $\lambda\gg0$, $S_\lambda$ 也具有相应性质.
\end{proposition}
\begin{proof}
见 \cite[Tags 01Z6, 086Q 及 01Z5]{SP} 及 \cite[命题 C.$6$--$7$]{TT90}.
\end{proof}
\begin{definition}\label{src:B.3.6}
称概形态射的性质 $\mathcal P$ \newterm{在极限下下降}{descends under limits}, 如果对每个拟紧且拟分离概形逆系统的极限 $S=\varprojlim_{\lambda\in\Lambda}S_\lambda$, 只要过渡映射仿射, 就有如下结论: 对每个指标 $0\in\Lambda$ 和每个拟紧且拟分离概形之间的态射 $g_0:X_0\to Y_0$, 其在 $S_\lambda$ 上的基变换为 $g_\lambda:X_\lambda\to Y_\lambda$, 在 $S$ 上的基变换为 $g:X\to Y$, 若 $g$ 具有 $\mathcal P$, 则对 $\lambda\gg0$, $g_\lambda$ 也具有 $\mathcal P$.
\end{definition}
\begin{proposition}[极限下态射性质的下降]\label{src:B.3.7}
概形态射的下列性质在极限下下降: 同构, 闭浸入, 开浸入, 仿射, 拟仿射, 有限, 拟有限, 固有, 射影, 拟射影, 分离, 单态射, 满射, 平坦, 局部有限呈示, 非分歧, 平展, 光滑, syntomic, 以及对某个固定整数 $d$, 每条纤维均连通且纯 $d$ 维的性质.
\end{proposition}
\begin{proof}
见 \cite[IV $8.10.5$]{EGA} 及 \cite[Tags 081C 及 05M5]{SP}.
\end{proof}

\subsection{展布与其他应用}\label{src:B.3.3-subsection}
概形的极限方法使我们可以将定义在环的余极限 (例如局部环 $R_{\mathfrak p}$) 上的对象\newterm{展布}{spreading out} 为定义在有限型环 (例如局部化 $R_f$) 上的对象.
\begin{proposition}[展布]\label{src:B.3.8}
设 $R$ 为环, $\mathfrak p\subseteq R$ 为素理想. 若 $X\to\Spec R_{\mathfrak p}$ 为有限呈示态射, 则存在元素 $f\notin\mathfrak p$ 及有限呈示态射 $X'\to\Spec R_f$, 使得 $X\simeq X'\times_{R_f}R_{\mathfrak p}$.
\end{proposition}
\begin{proof}
这是对 $R_{\mathfrak p}=\varinjlim_{f\notin\mathfrak p}R_f$ 应用极限下态射的下降 (\bookref{B.3.3}) 的直接结果.
\end{proof}
例如, 若 $Y$ 为函数域为 $K(Y)$ 的整概形, 则每个定义在 $K(Y)$ 上的有限呈示概形都可以延伸为定义在非空开子概形上的概形. 同样, 定义在 \S\bookref{B.5.3-subsection} 中的 Hensel 化 $R^h_{\mathfrak p}$ (相应地, 严格 Hensel 化 $R^{\mathrm{sh}}_{\mathfrak p}$) 上的有限呈示概形, 可以展布为 $\mathfrak p$ 的某个平展邻域 (相应地, 剩余域平凡的平展邻域) $\Spec R'\to\Spec R$ 上的有限呈示概形 $X'$.

作为 Noether 逼近的另一个典型应用, 我们说明如何将任意曲线族的性质归结为 Noether 基底上曲线族的性质.
\begin{proposition}\label{src:B.3.9}
设 $S$ 为拟紧且拟分离概形 (例如仿射概形), $C\to S$ 为概形之间的固有, 平坦且有限呈示态射, 每条几何纤维的维数至多为 $1$. 则存在 Cartesian 图
\[
\begin{tikzcd}[column sep=3em,row sep=2em]
C\arrow[r]\arrow[d]\arrow[dr,phantom,"\square" description]&C'\arrow[d]\\
S\arrow[r]&S'
\end{tikzcd}
\]
其中 $S'$ 为 $\mathbb Z$ 上的有限型概形, $C'\to S'$ 为概形之间的固有平坦态射, 每条几何纤维的维数至多为 $1$. 此外, 若 $C\to S$ 光滑, 则还可使 $C'\to S'$ 光滑.
\end{proposition}
\begin{proof}
将 $S=\varprojlim_{\lambda\in\Lambda}S_\lambda$ 写成 $\mathbb Z$ 上有限型概形逆系统的极限 (命题 \bookref{B.3.1}). 由于 $C\to S$ 有限呈示, 存在指标 $0\in\Lambda$ 和有限呈示态射 $C_0\to S_0$, 使得 $C\simeq C_0\times_{S_0}S$ (命题 \bookref{B.3.3}). 对每个 $\lambda>0$, 可定义 $C_\lambda=C_0\times_{S_0}S_\lambda$, 并有 Cartesian 图
\[
\begin{tikzcd}[column sep=3em,row sep=2em]
C\arrow[r]\arrow[d]\arrow[dr,phantom,"\square" description]&C_\lambda\arrow[r]\arrow[d]\arrow[dr,phantom,"\square" description]&C_0\arrow[d]\\
S\arrow[r]&S_\lambda\arrow[r]&S_0.
\end{tikzcd}
\]
由于 $C\to S$ 平坦且固有, 纤维维数至多为 $1$ (相应地, 光滑), 存在 $\lambda_0\in\Lambda$, 使得对所有 $\lambda\geq\lambda_0$, $C_\lambda\to S_\lambda$ 也如此 (命题 \bookref{B.3.7}). 现在对某个 $\lambda\geq\lambda_0$, 取 $S'=S_\lambda$ 和 $C'=C_\lambda$ 即可.
\end{proof}
由此可见, 若我们能证明 Noether 概形之间态射 $C'\to S'$ 的某些性质, 且这些性质在基变换下稳定, 则它们对 $C\to S$ 也成立. 在引理 \bookref{6.2.25} 中, 我们将对结点曲线族证明相同的性质.
\par\clearpage
\endgroup
